\chapter{Fixed-Branch Approximate Likelihoods and Same-Scalar Gradients}
\label{ch:bf-highdim-fixed-branch-same-scalar}
\chaptermark{Fixed-branch likelihoods and gradients}
\label{ch:bf-highdim-hmc-research}
\label{eq:bf-hd-hmc-potential}
\label{prop:bf-hd-hmc-gradient-contract}

\section{Why Fixed-Branch Discipline Is Shared Across The Two Lanes}
\label{sec:bf-hd-fixed-branch-purpose}

Both method chapters in the restarted high-dimensional block face the same
structural problem.  The exact Bayesian target is a filtering law and normalizer
sequence, but the implemented approximation uses discrete numerical structure
that may change with parameters: cloud level, merge policy, factor branch,
coordinate domains, ranks, pivots, grids, fitting points, or map choices.  If
those structural choices move while the parameter moves, then the reported
derivative is not automatically the derivative of one declared scalar.

This chapter therefore treats fixed-branch discipline as a cross-method target
contract, not as an appendix to one particular algorithm.

\section{A Canonical Branch Record}
\label{sec:bf-hd-branch-record}

A branch record should be read as a three-ledger object,
\begin{equation}
  \mathcal H(B)=\{\mathcal H_{\rm fixed},\mathcal H_{\rm diff},\mathcal H_{\rm diag}\},
  \label{eq:bf-hd-branch-ledger}
\end{equation}
where:
\begin{itemize}[leftmargin=0.28in]
  \item \(\mathcal H_{\rm fixed}\) contains the structural choices that define the
  scalar on the saved branch,
  \item \(\mathcal H_{\rm diff}(\theta)\) contains the parameter-dependent numerical
  values recomputed on that branch,
  \item \(\mathcal H_{\rm diag}\) contains residuals, condition numbers, rank
  vectors, branch-failure markers, and other diagnostic quantities.
\end{itemize}
For the sparse-grid lane, the fixed ledger includes the cloud, one-dimensional
rule family, merge/order policy, factor family, thresholds, preprocessing
policy, and initial-condition policy.  For the retained-object lane, the fixed
ledger includes coordinate domains, basis families, fitting points, weights,
ranks, sweep order, shifts, stabilization choices, retained-object construction
policy, and the defensive mass / reference-density conventions that define the
same branch.

Two evaluations use the same branch if and only if their fixed ledgers match:
\begin{equation}
  B(\theta_a)=B(\theta_b)
  \quad\Longleftrightarrow\quad
  \mathcal H_{\rm fixed}(\theta_a)=\mathcal H_{\rm fixed}(\theta_b).
  \label{eq:bf-hd-branch-equality}
\end{equation}
The differentiable ledger is allowed to change with the parameter; otherwise the
same-scalar derivative would not be testing the actual fixed branch.

\section[Fixed-branch approximate likelihoods]{Fixed-Branch Approximate Likelihoods Across The Three Operational Layers}
\label{sec:bf-hd-fixed-branch-likelihoods}

The restarted high-dimensional part now separates the operational burdens that
the old compressed block used to mix together.  The deterministic sparse-grid
lane exports a declared fixed-cloud scalar from
Chapter~\ref{ch:bf-highdim-fixed-cloud-sgqf-validation}.  The retained-object
TT/KR lane exports a declared fixed-branch scalar from
Chapter~\ref{ch:bf-highdim-squared-tt-fixed-branch-likelihoods}.  The present
chapter sits above those lane-specific constructions and states the common
same-scalar derivative logic they must both obey.

The three operational layers are:
\begin{equation}
\begin{aligned}
  \text{exact target:}&\quad \ell_T(\theta)=\sum_{t=1}^T\log Z_t(\theta),\\
  \text{declared approximate scalar:}&\quad \widehat\ell_T(\theta;B),\\
  \text{same-scalar derivative requirement:}&\quad
  \nabla_\theta \widehat\ell_T(\theta;B)
  \text{ differentiates the same declared } \widehat\ell_T.
\end{aligned}
  \label{eq:bf-hd-fixed-branch-three-layer}
\end{equation}
For the sparse-grid lane, \(B\) contains the fixed cloud, merge/order policy,
factor branch, and lane-specific veto rules.  For the retained-object lane,
\(B\) contains the fixed domains, fitting points, ranks, sweeps, shifts,
stabilization choices, and stored-field conventions that define the TT/KR scalar.
The chapter therefore treats fixed-branch discipline as a cross-method target
contract, not as an appendix to one particular algorithm.

For the actual-SV family, the ordering of obligations is now explicit.  First
one declares the target-correct scalar.  For the transformed actual-SV problem
that scalar is the cumulative likelihood built from the exact transformed
normalizers,
\[
  \widehat\ell_T^{\rm sv}(\theta;B)
  =
  \sum_{t=1}^T \log \widehat Z_t^{\rm sv}(\theta;B),
  \qquad
  \widehat Z_t^{\rm sv}(\theta;B)\approx Z_t(\theta).
\]
Only after that declaration does the same-scalar derivative contract become
meaningful.  Direct-likelihood quadrature, fixed-cloud reweighting, and
Zhao--Cui comparisons are legitimate only when they approximate this same
normalizer sequence.  The previous actual-SV ``Lane~B'' framing violated that
ordering by differentiating a Gaussian-closure replacement scalar built from
predictive observation moments of an augmented variable.  That scalar was not a
second admissible target for the intended inference problem; it was the wrong
scalar.  A finite-difference or autodiff check on that wrong scalar can show only
internal consistency of the replacement construction, not correctness for the
actual-SV likelihood.

\section{Finite-Difference Contract And Branch-Validity Logic}
\label{sec:bf-hd-fixed-branch-fd-contract}

The finite-difference contract is
\begin{equation}
  D_i(h)=
  \frac{\widehat\ell_T(\theta_0+h e_i;B)-\widehat\ell_T(\theta_0-h e_i;B)}{2h},
  \label{eq:bf-hd-fixed-branch-fd}
\end{equation}
with both perturbed evaluations reusing the same fixed ledger while recomputing
all parameter-dependent numerical values through the same declared equations.

For the retained-object lane, this means the domains, points, ranks, shifts, and
sweep order are unchanged, but the fitted core values
\begin{equation}
  C_{t,k}(\theta_0+h e_i;B)
  \quad\text{and}\quad
  C_{t,k}(\theta_0-h e_i;B)
  \label{eq:bf-hd-ttkr-perturbed-cores}
\end{equation}
are recomputed from the same saved fitting rule.  They are not copied from the
base parameter.  A row is branch-valid only if the two perturbations share:
\begin{enumerate}[label=(\roman*), leftmargin=0.28in]
  \item the same saved fixed ledger, and
  \item the same realized acceptance/failure path through the declared branch
  logic.
\end{enumerate}
If either side changes the accepted/failure pattern, the row is branch-invalid
and is not evidence about the derivative of the same scalar.

\paragraph{Copied-core and branch-identity audit protocol.}
The operational diagnostic should promote these prose rules into explicit
indicators.  For a step ladder \(h_i\), define
\begin{equation}
  I_\pm(h_i)=\mathbf{1}\{\mathcal H_{\rm fixed}(\theta_0\pm h_i e_j)=\mathcal H_{\rm fixed}(\theta_0)\},
  \label{eq:bf-hd-ttkr-branch-indicator}
\end{equation}
which checks branch-identity equality of the frozen fields, and
\begin{equation}
  K_\pm(h_i)=\mathbf{1}\{C_{t,k}(\theta_0\pm h_i e_j;B)\ \text{was recomputed by the fixed fitting rule}\},
  \label{eq:bf-hd-ttkr-copied-core-indicator}
\end{equation}
which checks that perturbed cores were not copied from the base parameter.
Record a decreasing-window indicator
\begin{equation}
  W_i=\mathbf{1}\{e(h_i)>e(h_{i+1})>e(h_{i+2})\},
  \label{eq:bf-hd-ttkr-decreasing-window}
\end{equation}
and classify the diagnostic status as one of:
\begin{equation}
\text{status}=
\begin{cases}
\text{branch identity failure},&\min_{i,\pm} I_\pm(h_i)=0,\\
\text{copied-core failure},&\min_{i,\pm} K_\pm(h_i)=0,\\
\text{decreasing-window agreement},&\max_i W_i=1,\\
\text{inconclusive: no decreasing window},&\text{otherwise}.
\end{cases}
  \label{eq:bf-hd-ttkr-fd-status}
\end{equation}
This makes the same-branch contract auditable.  If a frozen field differs, the
centered difference is comparing different functions.  If a differentiable field
is copied rather than recomputed, the centered difference is not evaluating the
fixed fitting rule.

\section[Value and derivative paths]{Value Path, Derivative Path, And What Is Actually Differentiated}
\label{sec:bf-hd-derivative-ledger}

The implementation burden of a same-scalar chapter is not just to state that a
gradient exists.  It must say which objects belong to the value path and which
objects belong to the derivative path.  A generic fixed-branch workflow has the
shape
\begin{equation}
\begin{aligned}
  &\text{value path:}
    &&\text{frozen structure} \to \widehat Z_t \to \widehat\ell_T,\\
  &\text{derivative path:}
    &&\text{same frozen structure} \to \dot{\widehat Z}_t \to
      \partial_i\widehat\ell_T.
\end{aligned}
  \label{eq:bf-hd-value-derivative-ledger}
\end{equation}
For the sparse-grid lane this means differentiating the same fixed cloud and the
same branch-conditioned Gaussian value path.  For the retained-object lane it
means differentiating the same retained-object fitting equations, mass
contractions, quotient updates, and carried-filter update equations without
changing domains, points, ranks, shifts, or sweep order.

The most important implementation distinction is between:
\begin{itemize}[leftmargin=0.28in]
  \item objects needed to close the current scalar derivative, and
  \item objects needed only to propagate the branch-consistent state to the next
  step.
\end{itemize}
This distinction is what keeps a same-scalar derivative from being confused with
a numerically similar but structurally different adaptive algorithm.

\paragraph{Retained-object derivative mechanics on the frozen branch.}
For the TT/KR lane, the forward and derivative paths should be read as the same
finite sequence of fixed equations with different outputs.  A useful teaching
layer is:
\begin{equation}
\begin{array}{p{0.41\linewidth}p{0.45\linewidth}}
\toprule
Forward scalar path & Derivative path\\
\midrule
\(\widetilde q_{t,j}=\widehat p_{t-1,j} f_j g_j J_j\) &
\(\dot{\widetilde q}_{t,j}=J_j(\dot{\widehat p}_{t-1,j}f_jg_j+\widehat p_{t-1,j}\dot f_jg_j+\widehat p_{t-1,j}f_j\dot g_j)\)\\
\(y_j=e^{c_t/2}\sqrt{\widetilde q_{t,j}}\) &
\(\dot y_j=e^{c_t/2}\dot{\widetilde q}_{t,j}/(2\sqrt{\widetilde q}_{t,j})\)\\
\(A_k\to N_k=A_k^\top W A_k+\rho I,\ d_k=A_k^\top Wy\) &
\(\dot A_k\to \dot N_k=\dot A_k^\top WA_k+A_k^\top W\dot A_k,\ \dot d_k=\dot A_k^\top Wy+A_k^\top W\dot y\)\\
\(N_k g_k=d_k\) &
\(N_k\dot g_k=\dot d_k-\dot N_k g_k\)\\
\(g_k\to C_k\to M_{>k}\to \widehat Z_t\) &
\(\dot g_k\to \dot C_k\to \dot M_{>k}\to \dot{\widehat Z}_t\)\\
\(\widehat p_t=a_t/\widehat Z_t\) &
\(\dot{\widehat p}_t=\dot a_t/\widehat Z_t-a_t\dot{\widehat Z}_t/\widehat Z_t^2\)\\
\bottomrule
\end{array}
  \label{eq:bf-hd-ttkr-derivative-teaching-table}
\end{equation}
The fixed linear solve is the key finite-dimensional derivative step:
\begin{equation}
  N_k g_k=d_k
  \quad\Longrightarrow\quad
  N_k\dot g_k=\dot d_k-\dot N_k g_k.
  \label{eq:bf-hd-ttkr-fixed-solve-derivative}
\end{equation}
Computation should solve \eqref{eq:bf-hd-ttkr-fixed-solve-derivative} with the
same numerical linear algebra used in the forward pass; it should not form an
explicit inverse.  After the core derivative is obtained, the remaining objects
follow by repeated product-rule and quotient-rule propagation through the same
frozen branch.

\paragraph{Mass-contraction and normalizer block.}
The squared-TT normalizer is not a separate conceptual object from the fitted
cores; it is a contraction of those same cores against the frozen basis mass
matrices.  Let the right mass contraction satisfy a recursion of the form
\begin{equation}
  M_{>j-1}[a,a']=
  \sum_{b,b',\ell,\ell'}
  C_j[a,\ell,b]
  M_{>j}[b,b']
  C_j[a',\ell',b']
  B_j[\ell,\ell'],
  \qquad j=D,\ldots,m+1,
  \label{eq:bf-hd-ttkr-mass-recursion}
\end{equation}
for a retained prefix of dimension \(m\), with basis mass matrix \(B_j\) frozen
by the branch.  Continuing the same recursion through \(j=1\) gives the full
square mass.  Differentiate the suffix recursion by the product rule:
\begin{equation}
\begin{aligned}
  \dot M_{>j-1}[a,a']
  &=
  \sum_{b,b',\ell,\ell'}
  \dot C_j[a,\ell,b]
  M_{>j}[b,b']
  C_j[a',\ell',b']
  B_j[\ell,\ell']\\
  &\quad+
  \sum_{b,b',\ell,\ell'}
  C_j[a,\ell,b]
  \dot M_{>j}[b,b']
  C_j[a',\ell',b']
  B_j[\ell,\ell']\\
  &\quad+
  \sum_{b,b',\ell,\ell'}
  C_j[a,\ell,b]
  M_{>j}[b,b']
  \dot C_j[a',\ell',b']
  B_j[\ell,\ell'],
  \qquad j=D,\ldots,m+1.
\end{aligned}
  \label{eq:bf-hd-ttkr-dot-mass-recursion}
\end{equation}
Starting from the terminal contraction, this recursion produces the square mass
\(R_t\) and its derivative \(\dot R_t\).  Assume here that the fixed defensive
density is normalized under the declared full reference measure,
\(\int\lambda_t\dd\mu_{1:D}=1\).  If
\begin{equation}
  \zeta_t=R_t+\tau_t,
  \qquad
  \widehat Z_t=e^{-c_t}\zeta_t,
  \label{eq:bf-hd-ttkr-normalizer-block}
\end{equation}
then, with \(c_t\) and \(\tau_t\) frozen on the branch,
\begin{equation}
  \dot{\widehat Z}_t=e^{-c_t}\dot R_t.
  \label{eq:bf-hd-ttkr-dot-normalizer}
\end{equation}
This is the mass-contraction block the derivative must preserve.  Omitting one of
the product-rule terms in \eqref{eq:bf-hd-ttkr-dot-mass-recursion} means the
reported gradient no longer differentiates the same scalar normalizer.

\paragraph{Carried-filter and next-step closure block.}
The retained numerator is another contraction of the same represented
adjacent-state object.  Write
\begin{equation}
  a_t(z_t;\beta)=
  \int \widehat q_t(z_t,z_{t-1};\beta)\dd z_{t-1},
  \qquad
  \dot a_t(z_t;\beta)=
  \int \dot{\widehat q}_t(z_t,z_{t-1};\beta)\dd z_{t-1}.
  \label{eq:bf-hd-ttkr-carried-numerator}
\end{equation}
The normalized carried filter then satisfies the quotient rule
\begin{equation}
  \dot{\widehat p}_t(z_t;\beta)=
  \frac{\dot a_t(z_t;\beta)}{\widehat Z_t(\beta)}
  -
  \frac{a_t(z_t;\beta)\dot{\widehat Z}_t(\beta)}{\widehat Z_t(\beta)^2}.
  \label{eq:bf-hd-ttkr-dot-carried-filter}
\end{equation}
This is the crucial bridge between time steps.  At time \(t+1\), the fixed-branch
target contains \(\widehat p_t\), so its derivative contains
\(\dot{\widehat p}_t\).  Without saving the carried-filter derivative, the
next-step gradient is incomplete even if the current-step score is numerically
correct.

\paragraph{Saved retained-filter object and query rule.}
The phrase “save the carried object” must denote a finite-dimensional evaluator,
not an informal function name.  For retained state dimension $m>1$, let the
retained block use the product basis
\begin{equation}
  b_t(z_t)=b_{t,1}(z_{t,1})\otimes\cdots\otimes b_{t,m}(z_{t,m})
  \in\R^{p_{\rm ret}},
  \qquad
  p_{\rm ret}=\prod_{j=1}^m p_j.
  \label{eq:bf-hd-ttkr-retained-product-basis}
\end{equation}
The link from the retained TT prefix to a coefficient matrix must be explicit.
Because the bases form a tensor-product basis, there is a matrix
\(A_t\in\R^{p_{\rm ret}\times R_m}\) satisfying
\begin{equation}
  G_m(z_t;\beta)=b_t(z_t)^\top A_t(\beta),
  \qquad
  \dot G_m(z_t;\beta)=b_t(z_t)^\top\dot A_t(\beta).
  \label{eq:bf-hd-ttkr-retained-prefix-coefficients}
\end{equation}
The current fixed branch also requires the retained defensive density to have a
declared coefficient matrix
\begin{equation}
  \lambda_{t,\mathrm{ret}}(z_t)
  =b_t(z_t)^\top\Lambda_t^{\rm ret}b_t(z_t),
  \qquad
  \dot\Lambda_t^{\rm ret}=0.
  \label{eq:bf-hd-ttkr-defensive-retained-coefficients}
\end{equation}
If the selected defensive density has no exact representation of this form,
the implementation must store it as a separate evaluator; it may not silently
drop the defensive contribution from a \(Q_t\)-only saved object.

With \(c_t,\tau_t,\Lambda_t^{\rm ret}\), the bases, and the reference measure
frozen, the coefficient matrices derived from the Chapter~36b contraction are
\begin{equation}
\begin{aligned}
  &Q_t,\dot Q_t,\Lambda_t^{\rm ret}
  \in\R^{p_{\rm ret}\times p_{\rm ret}},\\
  Q_t
  &=
  e^{-c_t}\left(
  A_tM_{>m}A_t^\top+\tau_t\Lambda_t^{\rm ret}\right),\\
  \dot Q_t
  &=
  e^{-c_t}\left(
  \dot A_tM_{>m}A_t^\top+
  A_t\dot M_{>m}A_t^\top+
  A_tM_{>m}\dot A_t^\top\right).
\end{aligned}
  \label{eq:bf-hd-ttkr-retained-Q-construction}
\end{equation}
Consequently, the retained numerator and its derivative use the matching
quadratic evaluators.  Here \(b_t(z_t)\in\R^{p_{\rm ret}}\) is a column vector,
so every product below is scalar:
\begin{equation}
  a_t(z_t;\beta)=b_t(z_t)^\top Q_t(\beta)b_t(z_t),
  \qquad
  \dot a_t(z_t;\beta)=b_t(z_t)^\top\dot Q_t(\beta)b_t(z_t),
  \label{eq:bf-hd-ttkr-retained-Q}
\end{equation}
Normalize the stored evaluator by
\begin{equation}
  P_t=\frac{Q_t}{\widehat Z_t},
  \qquad
  \dot P_t=
  \frac{\dot Q_t}{\widehat Z_t}
  -
  \frac{Q_t\dot{\widehat Z}_t}{\widehat Z_t^2}.
  \label{eq:bf-hd-ttkr-retained-P}
\end{equation}
For next-step query points $z_t^{\rm query}(j)$, build the retained-basis matrix
\begin{equation}
  B^{\rm query}[j,:]=b_t(z_t^{\rm query}(j))^\top,
  \qquad
  b_t(z_t^{\rm query}(j))^\top\in\R^{1\times p_{\rm ret}},
  \qquad B^{\rm query}\in\R^{M\times p_{\rm ret}},
  \label{eq:bf-hd-ttkr-retained-query-basis}
\end{equation}
and evaluate the saved carried object by
\begin{equation}
  \widehat p_t^{\rm ref}[j]
  =
  B^{\rm query}[j,:]P_tB^{\rm query}[j,:]^\top,
  \qquad
  \dot{\widehat p}_t^{\rm ref}[j]
  =
  B^{\rm query}[j,:]\dot P_tB^{\rm query}[j,:]^\top.
  \label{eq:bf-hd-ttkr-retained-query-rule}
\end{equation}
The next-step query rule must also state the coordinate system and Jacobian
ownership explicitly.  In the current fixed-adjacent route, the stored carried
marginal is a density under the retained-block \emph{reference} measure, and the
saved evaluator therefore consumes retained reference coordinates directly.
Physical fitting points are converted to reference coordinates by the route's
coordinate map before the saved evaluator is queried.  The Jacobian of that
physical-to-reference conversion is owned by the step-(t+1) target-construction
path, not by the stored evaluator itself.  Thus the next-step query rule is that
\(z_t^{\rm query}(j)\) denotes the previous-state block of the fixed branch's
fitting point \emph{after conversion into the retained reference coordinates on
which \eqref{eq:bf-hd-ttkr-retained-query-rule} is defined}.  A later
implementation may store \(Q_t\) in TT, low-rank, or sparse form, but only if it
preserves the same evaluator and derivative evaluator in
\eqref{eq:bf-hd-ttkr-retained-query-rule}.  This is the stored-field convention
that closes the derivative recursion on the same branch.

\paragraph{Route-audit note.}
The release-oriented text above is a contract statement rather than a source or
code audit.  A final promotion should verify that the actual fixed-branch route
uses the same reference measure and conversion ownership before this paragraph is
promoted beyond a documented convention.

\paragraph{Branch-manifest reporting hook.}
Derivative and finite-difference reports for this lane should therefore carry a
TT/KR branch-manifest identifier together with the saved-field convention, so
that branch-valid rows certify equality of the actual retained evaluator rather
than only superficial notation.

\paragraph{Compact derivative-reading checklists.}
A later implementation review should be able to verify the derivative lane block
by block:
\begin{equation}
\begin{array}{ll}
\text{squared-density block:} & \overline q_t=e^{-c_t}(\phi_t^2+\tau_t\lambda_t),\ \partial(\phi_t^2)=2\phi_t\dot\phi_t;\\[1mm]
\text{mass-contraction block:} & R_t=\int \phi_t^2,\ \dot R_t\ \text{from product-rule contractions};\\[1mm]
\text{fixed-solve block:} & N_kg_k=d_k\ \mapsto\ N_k\dot g_k=\dot d_k-\dot N_kg_k;\\[1mm]
\text{carried-filter block:} & \widehat p_t=a_t/\widehat Z_t,\ \dot{\widehat p}_t=\dot a_t/\widehat Z_t-a_t\dot{\widehat Z}_t/\widehat Z_t^2.
\end{array}
  \label{eq:bf-hd-ttkr-derivative-checklists}
\end{equation}
The corresponding failure-mode summaries are:
\begin{itemize}[leftmargin=0.28in]
  \item \textbf{squared-density failure:} differentiating a different fitted
  square root than the one actually squared,
  \item \textbf{mass-contraction failure:} omitting a product-rule term in the
  contraction recursion,
  \item \textbf{fixed-solve failure:} changing the normal equations or copying a
  stale core at perturbed parameters,
  \item \textbf{carried-filter failure:} feeding time \(t+1\) without the saved
  dotted retained object or with an undeclared query rule.
\end{itemize}
These lists are not replacement derivations; they are compact audit objects that
make it harder for a future implementation to silently drift away from the
mathematical branch defined here.

\section{A numerical score from a conditional smoothing proposal}
\label{sec:bf-hd-quadratic-reference}
\sectionmark{Numerical score from a smoothing proposal}

An independent likelihood calculation gives a useful way to assess the score of
an approximate filter.  The comparison must hold the observations, parameter
coordinates, initial law, transition law and observation law fixed.  For the
observed sequence $y_{1:T}$, write the complete latent trajectory as
$x=x_{0:T}$ and define
\begin{equation}
  p_\theta(x,y_{1:T})
  =\mu_\theta(x_0)\prod_{t=1}^T
       f_\theta(x_t\mid x_{t-1})g_\theta(y_t\mid x_t),
  \qquad
  Z(\theta)=\int p_\theta(x,y_{1:T})\,dx.
  \label{eq:bf-hd-reference-path-target}
\end{equation}
The quantity to estimate is the likelihood score
$s(\theta_0)=\nabla_\theta\log Z(\theta_0)$.  It contains no parameter-prior
term.  A Gaussian observation law, for example, contributes both its quadratic
residual and its log-determinant normalization; changing an observation variance
requires recomputing both.  Omitting the normalization would differentiate a
different likelihood.

The filtering and backward-smoothing constructions of
\citet[Algorithms 1, 2 and 4]{zhao2024ttsequential} can provide a proposal for
$x$ at a fixed parameter $\theta_0$.  First reproduce their filtering and
smoothing experiment, including its model equations and reported quantities.
Then assess whether a proposal constructed for our specified data and parameter
supports likelihood integration.  The local quadratic score calculation below
is a separate extension; it is not a score algorithm asserted by that paper.
In particular, paths drawn from the predator--prey \emph{joint} posterior over
parameters and states cannot simply be relabeled as paths from the conditional
smoother at one chosen parameter.  The latter requires its own proposal.

Let $q_0(x)$ be that conditional path proposal, fitted once at $\theta_0$, and
draw $N$ paths $X_1,\ldots,X_N$ from it.  For the importance identity to cover
the intended likelihood, $q_0(x)$ must be positive wherever
$p_\theta(x,y_{1:T})$ is positive throughout the neighborhood being evaluated.
Keeping the proposal and paths fixed gives
\begin{align}
  Z(\theta)
    &=\mathbb E_{q_0}\!\left[
        \frac{p_\theta(X,y_{1:T})}{q_0(X)}\right],
       \label{eq:bf-hd-reference-importance-identity}\\
  \widehat Z_N(\theta)
    &=\frac1N\sum_{i=1}^N \exp a_i(\theta),
  \qquad
  a_i(\theta)=\log p_\theta(X_i,y_{1:T})-\log q_0(X_i).
       \label{eq:bf-hd-reference-importance-estimate}
\end{align}
Thus every new parameter point needs new joint-density evaluations, but it does
not need a new fitted tensor train.  This does not copy fitted cores while
claiming to differentiate the adaptive fitted likelihood discussed earlier.
It defines a different, explicit finite-sample scalar,
$\log\widehat Z_N(\theta)$, whose proposal is constant and whose entire target
density is recomputed.  Under ordinary integrability conditions
$\widehat Z_N$ is an unbiased estimate of $Z$ when the paths have law $q_0$;
its logarithm and its score are generally biased at finite $N$.

To avoid subtracting two large log likelihoods, normalize the baseline weights:
\begin{equation}
  \alpha_i=
  \frac{\exp a_i(\theta_0)}{\sum_{k=1}^N\exp a_k(\theta_0)},
  \qquad
  \delta_i(\theta)=
  \log p_\theta(X_i,y_{1:T})-\log p_{\theta_0}(X_i,y_{1:T}).
  \label{eq:bf-hd-reference-baseline-weights}
\end{equation}
Substituting $\exp a_i(\theta)=\exp a_i(\theta_0)\exp\delta_i(\theta)$
into the numerator of the likelihood ratio gives
\begin{align}
  \frac{\widehat Z_N(\theta)}{\widehat Z_N(\theta_0)}
   &=\frac{\sum_i\exp a_i(\theta_0)\exp\delta_i(\theta)}
           {\sum_i\exp a_i(\theta_0)}
     =\sum_i\alpha_i\exp\delta_i(\theta),\\
  \widehat r_N(\theta)
   &=\log\widehat Z_N(\theta)-\log\widehat Z_N(\theta_0)
     =\log\sum_i\alpha_i\exp\delta_i(\theta).
  \label{eq:bf-hd-reference-log-ratio}
\end{align}
Evaluate these sums by log-sum-exp.  A parameter-independent multiplicative
constant in the recorded proposal density cancels in the ratio, but prevents
recovery of an absolute likelihood unless that constant is restored.  A path
with zero target weight remains in the original sample count.  When a recorded
log density is $-\infty$, the direct log-weight expression in
\eqref{eq:bf-hd-reference-importance-estimate} avoids the undefined subtraction
$(-\infty)-(-\infty)$.

Choose positive parameter scales $d_1,\ldots,d_p$, set
$D=\operatorname{diag}(d_1,\ldots,d_p)$, and express nearby parameters as
$\theta=\theta_0+hDz$, where $h>0$ is a dimensionless radius.  This lets, for
example, a population-capacity parameter and a growth-rate parameter receive
comparable relative perturbations despite their different units.  At many
specified points $z_j$, fit the full quadratic
\begin{equation}
  (\widehat a,\widehat b,\widehat C)
  =\arg\min_{a,b,C=C^\mathsf T}
  \sum_{j=1}^M\left[
       \widehat r_N(\theta_0+hDz_j)
       -a-b^\mathsf Tz_j-\tfrac12z_j^\mathsf TCz_j
       \right]^2.
  \label{eq:bf-hd-reference-quadratic-fit}
\end{equation}
The intercept absorbs fitting error rather than forcing it into the slope.
All cross terms are present: the design has
$1+p+p(p+1)/2$ columns, hence 28 columns for six parameters and 10 for three.
Use a full-rank design and report its condition number; a ridge would change the
estimated slope and requires a separate justification.

The chain rule supplies the units of the answer.  Since
$z=h^{-1}D^{-1}(\theta-\theta_0)$, the fitted score and Hessian at the center are
\begin{equation}
  \widehat s_h(\theta_0)=h^{-1}D^{-\mathsf T}\widehat b,
  \qquad
  \widehat H_h(\theta_0)
     =h^{-2}D^{-\mathsf T}\widehat C D^{-1}.
  \label{eq:bf-hd-reference-quadratic-derivatives}
\end{equation}
For a diagonal $D$, the $k$th score component is $\widehat b_k/(hd_k)$.
If the fitted coordinates are log parameters, this is the score with respect to
those log parameters; conversion to physical parameters requires their
Jacobian.  It does not follow from changing the labels of the reported values.

A scalar Gaussian example makes the scaling concrete.  With
$y\sim\mathcal N(\theta,v)$, the exact log-likelihood difference is
\begin{equation}
  \ell(\theta_0+hdz)-\ell(\theta_0)
   =hd\frac{y-\theta_0}{v}z-\frac{h^2d^2}{2v}z^2.
  \label{eq:bf-hd-reference-gaussian-example}
\end{equation}
The fit therefore returns $b=hd(y-\theta_0)/v$ and $C=-h^2d^2/v$;
dividing by $hd$ and $h^2d^2$ recovers the exact score and curvature.
For $y=2$, $\theta_0=1$, $v=4$, $d=2$ and $h=0.1$, the coefficients are
$b=0.05$ and $C=-0.01$, while the score is $0.25$ and the curvature is $-0.25$.
This exact recovery tests the regression and coordinate conversion.  It says
nothing about whether a numerical state-space likelihood supplied to the same
fit is accurate.

Use antithetic design points: whenever $z$ is present, include $-z$ with the
same weight.  The odd columns $z_k$ are then orthogonal to the even intercept
and quadratic columns.  A third-order Taylor term still contributes to the
fitted slope: for a sufficiently smooth log likelihood and a fixed,
well-conditioned scaled design, $\widehat b=hD^\mathsf Ts+O(h^3)$ and the
resulting truncation error in the score is $O(h^2)$.  Symmetry removes leakage
from even terms into the slope; it does not remove cubic bias.  With the same
symmetric design, adding quadratic terms to a linear fit leaves its slope
unchanged in exact arithmetic.  The quadratic fit additionally describes
curvature and improves the representation of nearby likelihood values.
Consequently agreement between those two slopes is a design check, not
independent evidence of score accuracy.

The complete calculation can be organized as follows.
\begin{enumerate}[label=(\arabic*),leftmargin=0.28in]
  \item Fix the data, parameter coordinates and full probability laws.  Check
  the author experiment separately, then fit conditional proposals at the
  required $\theta_0$ and save paths and their proposal densities.
  \item Recompute every initial, transition and observation log density at
  $\theta_0$.  Form the baseline weights and absolute likelihood when proposal
  normalization is known.  Compare that likelihood with an independent
  calculation for the same data and laws.
  \item Evaluate the common-path likelihood differences at a symmetric
  space-filling design, heldout points and positive/negative coordinate probes.
  A diagnostic starting design uses 1024 fitted points and 256 heldout points
  at radii $0.04$, $0.02$ and $0.01$; these are hypotheses to check for the
  particular target, not universal numerical settings.
  \item Fit the full quadratic at each radius, transform its slope back to
  parameter units, and retain heldout residuals, condition numbers, central
  differences, normalized weights and effective sample sizes.
  \item Repeat with independent fitted proposals and independent path draws,
  and assess sensitivity to rank, sample size and radius.  Report the score
  coordinates together with their observed numerical variation and uncertainty.
\end{enumerate}

The weight diagnostics must cover every nearby parameter, not just the center.
For normalized weights $\beta_i(\theta)$ proportional to
$\exp a_i(\theta)$, report
\begin{equation}
  \operatorname{ESS}(\theta)=\frac{1}{\sum_i\beta_i(\theta)^2},
  \qquad \max_i\beta_i(\theta).
  \label{eq:bf-hd-reference-local-ess}
\end{equation}
These reveal domination by a few sampled paths.  They cannot reveal a region
that the proposal never samples.  If $q_0$ has support $A$ smaller than the
model's support, the importance integral converges to
$\int_Ap_\theta(x,y)\,dx$, and its derivative can differ from the full score
even with large ESS.  Rank and domain checks, independently fitted proposals,
and an independent likelihood reference are needed to investigate that error.

Finally, the likelihood values at different design points reuse the same
paths, so their Monte Carlo errors are correlated.  Ordinary regression
standard errors do not measure uncertainty in the likelihood score.  Repeating
only backward sampling from one fitted tensor train measures variability
conditional on that fit.  Repeating the fit as well explores fitting
variability, but neither operation by itself removes a shared support error.
The conditional sampling contribution can be estimated without evaluating
any additional likelihoods.  Divide the $N$ independent proposal paths into
$G$ equal groups, remove each group in turn, recompute the log-likelihood
ratios on the retained paths, and refit the same quadratic.  Write
$\widehat s_{(-g),k}$ for coordinate $k$ of the score after removing group $g$,
and $\overline s_{(-\cdot),k}=G^{-1}\sum_g\widehat s_{(-g),k}$.  The
corresponding jackknife standard error is
\begin{equation}
  \widehat{\operatorname{se}}_{\mathrm{path},k}
  =\left[\frac{G-1}{G}\sum_{g=1}^{G}
    \bigl(\widehat s_{(-g),k}-\overline s_{(-\cdot),k}\bigr)^2
    \right]^{1/2}.
  \label{eq:bf-hd-reference-jackknife}
\end{equation}
Removing the same group at every parameter point preserves the dependence
that matters for the fitted slope.  Twenty groups provide an initial
conditional diagnostic for $N=10{,}000$ paths; they do not provide twenty
independent fitted proposals.  The estimate is unreliable when a few paths
carry nearly all weight, and is unavailable when deleting a group leaves
zero total weight.  Report the mass carried by each group and compare the
score from the first $N/4$, $N/2$ and all $N$ paths as additional sampling
checks.  None of these calculations measures missing support or quadratic
truncation error.

A small heldout residual establishes that a quadratic represents the computed
local likelihood well.  Calling its derivative an oracle additionally requires
same-target likelihood agreement and score stability under the measured
sampling, fitting, rank and radius changes, with uncertainty reported in each
parameter coordinate.

\subsection{What the published experiments establish}
\label{sec:bf-hd-reference-replication}

The experiments in Sections 6.3--6.4 of
\citet{zhao2024ttsequential} assess how closely the tensor-train smoothing
proposal approximates a joint path posterior.  Their importance-weight ESS
is a useful test of that approximation.  It is not a direct test of a marginal
likelihood derivative.  To establish the reference calculation above, we first
ran the original fitting and backward-smoothing operations on the SIR model
with eighteen states and on the predator--prey model, then examined likelihoods
and numerical derivatives on our separately specified fixed-parameter targets.

The publication comparison used twenty observations, 33 polynomial degrees of
freedom and five ALS sweeps, with SIR rank caps 10, 20 and 40 and predator--prey
rank cap 20.  Every smoothing draw contained 10,000 paths.  The original MATLAB
data, fitted tensors and figure arrays were unavailable, so the computation
used a recorded Octave realization.  Table~\ref{tab:bf-hd-paper-ess-replication}
compares the local measurements with values extracted from the vector figures.
The predator--prey figure has ten observation updates; its horizontal axis
must not be stretched to twenty.  The paper's separate statement of
approximately 40\% ESS after twenty observations is a prose comparison.

\begin{table}[htbp]
\centering
\small
\begin{tabular}{llrrrr}
\toprule
Model & Proposal & $T$ & Paper median & Local ESS & Draws\\
\midrule
Predator--prey & Linear & 10 & 31.24\% & 88.39\% & 1\\
Predator--prey & Nonlinear & 10 & 80.71\% & 81.91\% & 40\\
SIR & Rank 10 & 20 & 15.40\% & 11.30\% & 40\\
SIR & Rank 20 & 20 & 57.84\% & 36.55\% & 40\\
SIR & Rank 40 & 20 & 69.34\% & 50.08\% & 1\\
\bottomrule
\end{tabular}
\caption{Joint-path ESS as a percentage of 10,000 samples.  Local rows with
40 draws report the median conditional on one fitted proposal and one dataset;
they are not forty independent repetitions of the paper's full experiment.
Single draws and these conditional medians have different sampling uncertainty
from the published medians.}
\label{tab:bf-hd-paper-ess-replication}
\end{table}

The nonlinear predator--prey measurement at $T=10$ is close to the published
median, but the much larger local linear ESS means that the published
nonlinear advantage is not reproduced on this realization.  At $T=20$,
the local linear and nonlinear conditional medians are 84.40\% and 75.18\%,
with interquartile ranges 80.19--86.45\% and 72.01--77.12\%, respectively.
The SIR measurements retain the expected ordering with rank, while their
levels differ from the published medians.  These comparisons are descriptive:
they neither establish a statistical ranking nor refute the paper.

The released demonstration scripts also differ from the printed experiment
in the RK4 fourth stage, initial-state generation, some predator--prey
parameters and priors, and the number of subsequent ALS sweeps.  Running those
settings separately gives terminal predator--prey ESS values of 33.84\% for
linear and 60.29\% for nonlinear preconditioning, and SIR rank-40 ESS of
25.49\%, each from one draw.  The qualitative predator--prey advantage appears
there.  Because several settings and the simulated data change together, this
does not identify which difference explains the discrepancy.  It shows why
an original-code comparison and a printed-experiment comparison must retain
their respective model definitions.

There is also a precise distinction between the source's returned diagnostic
and the marginal log likelihood needed here.  The released
\texttt{full\_sol.smooth} method returns the mean of the finite log importance
weights under the name \texttt{lml}.\footnote{Audited author commit
\texttt{80034dccb99eb1d86284a1839b4a12067d13b9da},
\href{https://github.com/DeepTransport/tensor-ssm-paper-demo/blob/80034dccb99eb1d86284a1839b4a12067d13b9da/models/full_sol.m\#L188-L197}{\texttt{models/full\_sol.m}, lines 188--197}.}
For finite $a_n=\log p_\theta(x^{(n)},y)-\log q_0(x^{(n)})$, the distinction is
\begin{equation}
  \frac1N\sum_{n=1}^N a_n
  \quad\hbox{versus}\quad
  \log\left(\frac1N\sum_{n=1}^N e^{a_n}\right).
  \label{eq:bf-hd-source-logweight-distinction}
\end{equation}
These quantities differ unless all weights are equal.  If the relevant
expectations exist, the population version of the first is
$\log Z(\theta)-\mathrm{KL}(q_0\Vert p_\theta(\cdot\mid y))$;
it therefore measures a different quantity.  Dropping zero-weight samples
would introduce a further conditioning change.  Our likelihood estimate uses
the second expression with all $N$ paths, including valid zero weights,
and the complete cumulative proposal density.  The paper's ESS experiments
do not depend on the returned mean-log-weight diagnostic, so correcting its
use as a likelihood does not establish an error in those published results.


\subsection{Actual fixed-target likelihoods and numerical scores}
\label{sec:bf-hd-final-score-values}

The likelihood calculation and the score calculation use the same finite
importance-sampling value program.  At $T=20$ the three rank-20 fitted
proposals give the following log-likelihood values at the fixed target.  The
bootstrap column is an independent four-replication calculation with
$524{,}288$ paths; its MCSE is the standard error across those replications.

\begin{table}[htbp]
\centering
\scriptsize
\begin{tabular}{lrrrrr}
\toprule
Model & seed 2 & seed 17 & seed 29 & bootstrap & MCSE\\
\midrule
Predator--prey & $-97.344056298$ & $-97.344035221$ & $-97.344030690$ & $-97.342973851$ & $0.006574815$\\
SIR $d=18$ & $-678.067711876$ & $-678.045207144$ & $-678.052340241$ & $-678.077466314$ & $0.014864040$\\
\bottomrule
\end{tabular}
\caption{Finite-target log likelihoods from the conditional rank-20 fits and
an independent bootstrap reference.  These are actual values of the same
observation and model laws.  The bootstrap has finite-particle bias that is
not estimated here.}
\label{tab:bf-hd-final-likelihood-values}
\end{table}

The local derivative is obtained by fitting the full quadratic
$\widehat\ell(z)=a+b^\top z+\frac12z^\top C z$ to the log-likelihood values
at common paths and returning $b/(hD)$.  The table reports the smallest
radius in each completed high-precision series.  The between-fit column is a
Student-$t$ half-width over three independently fitted proposals when three
are available; the path column is the equal-group conditional jackknife
standard error.  These intervals omit shared-support error, finite-particle
bias, and radius bias.

\begin{table}[htbp]
\centering
\scriptsize
\begin{tabular}{llrrrr}
\toprule
Case & $h$ & coordinate & score & between-fit half-width & path SE\\
\midrule
PP, $N=10^4$ & .00125 & $r$ & $-33.304023$ & .398694 & .114709\\
 & & $K$ & $-.55864759$ & .0358593 & .00589052\\
 & & $a$ & $.020413745$ & .000206453 & .000117006\\
 & & $s$ & $4.5071742$ & .0768389 & .0207010\\
 & & $u$ & $-8.7160499$ & .0141339 & .0237740\\
 & & $v$ & $10.847440$ & .0118746 & .0296156\\
\midrule
PP, $N=10^5$ & .0003125 & $r$ & $-33.051124$ & .187941 & .0330587\\
 & & $K$ & $-.55078735$ & .0147903 & .00174556\\
 & & $a$ & $.020627942$ & .000094596 & .0000323612\\
 & & $s$ & $4.4686647$ & .0283026 & .00586542\\
 & & $u$ & $-8.7462792$ & .0265643 & .00661567\\
 & & $v$ & $10.883935$ & .0280803 & .00828419\\
\midrule
SIR rank 40 & .00015625 & $\log\kappa$ & $112.685217$ & --- & $12.4285$\\
 & & $\log\nu$ & $-65.165591$ & --- & $5.4720$\\
 & & $\log R$ & $5.576446$ & --- & $.08618$\\
\bottomrule
\end{tabular}
\caption{Actual full-quadratic score estimates on the fixed $T=20$ targets.
The PP $10^5$ run has minimum ESS $99{,}892$ and maximum held-out RMS
$4.60\times10^{-8}$; the PP $10^4$ run has minimum ESS $9{,}824$ and RMS
$2.95\times10^{-6}$.  The SIR rank-40 run has ESS $4{,}549$ and RMS
$7.24\times10^{-7}$, but only one fitted proposal, so an independent-fit
interval is unavailable.}
\label{tab:bf-hd-final-score-values}
\end{table}

For PP, increasing the path count moves the score toward the independent
bootstrap vector
$(-32.905236,-.543943,.02052955,4.487344,-8.732397,10.865452)$ and reduces
the conditional path error.  For SIR, shrinking the neighborhood repairs the
rank-20 wide-radius overlap failure: the wide run has minimum ESS only
$1.17$--$1.34$ and held-out RMS $0.766$--$0.846$, whereas the rank-40
small-radius run has ESS $4{,}549$ and RMS $7.24\times10^{-7}$.  The rank-20
small-radius three-fit run reached its declared time limit after eight of nine
estimates, so it is incomplete and is not aggregated.  These results support a
well-resolved numerical derivative for the finite value program, with explicit
sampling and proposal uncertainty; they do not establish an exact oracle or a
filter-ranking result.


\section[Marginal importance correction]{Marginal importance correction across ancestor components}
\label{sec:bf-hd-marginal-importance-correction}

A particle moved by LEDH may lie far from the transition distribution of its
own ancestor and still be plausible under another ancestor.  Evaluating only
its generating component can then give a very different weight from evaluating
the complete predictive mixture.  Both constructions have an importance-sampling
interpretation.  The distinction concerns the conditioning variable retained in
the weight, rather than a change in the flow differential equation.

Condition on the previous weighted cloud.  Let $\omega_j>0$ denote its normalized
weights, $f_j(x)$ the Gaussian transition density from ancestor $j$, and
$q_j(x)$ the density after its invertible affine flow.  Write the transition
mean as $m_j$, its covariance as $Q$, and its flow as $F_j(z)=B_jz+b_j$.
The matrices and offsets are fixed with respect to the newly sampled innovation,
although they depend on the previous cloud, the observation and the parameter.
The change-of-variables formula gives
\begin{equation}
 q_j(x)=\mathcal N(x;B_jm_j+b_j,B_jQB_j^\top),\qquad
 \bar f(x)=\sum_j\omega_j f_j(x),\qquad
 \bar q(x)=\sum_j\omega_j q_j(x).
 \label{eq:bf-hd-marginal-component-densities}
\end{equation}
The positive-definite covariance and invertibility assumptions matter: a
singular transition or an innovation-dependent nonlinear map does not have the
Gaussian proposal density in this equation.

Suppose one particle $X_i$ is sampled from each $q_i$, with observation likelihood
$g(y\mid X_i)$.  The two unnormalized weights are
\begin{align}
 a_i^{\rm ancestor}
 &=\omega_i g(y\mid X_i)\frac{f_i(X_i)}{q_i(X_i)},\\
 a_i^{\rm mixture}
 &=\omega_i g(y\mid X_i)\frac{\bar f(X_i)}{\bar q(X_i)}.
 \label{eq:bf-hd-marginal-two-weights}
\end{align}
The outer $\omega_i$ remains in the mixture construction because the design
still has one draw from each component.  The inner mixture weights describe
which ancestors explain the evaluated state.  They are the same incoming
weights, used for a different purpose.  For any integrable test function $h$,
\begin{align}
 \mathbb E\!\left[\sum_i a_i^{\rm mixture}h(X_i)\right]
 &=\sum_i\omega_i\int q_i(x)g(y\mid x)
       \frac{\bar f(x)}{\bar q(x)}h(x)\,dx\\
 &=\int g(y\mid x)\bar f(x)h(x)\,dx.
 \label{eq:bf-hd-marginal-weight-identity}
\end{align}
The cancellation uses $\sum_i\omega_iq_i=\bar q$.  The ancestor construction
has the same expectation after cancelling each $q_i$ separately.  Thus the
mixture change is not a proof that the ancestor importance identity was wrong.
It changes the realized weights and can change variance substantially.  Neither
identity proves that a subsequent deterministic moment reset preserves the
exact filtering distribution.

The singleton limit also explains how one implementation serves both choices.
For the ancestral selection $S_i=\{i\}$, the ratio
$\sum_{j\in S_i}\omega_j f_j(X_i)/
 \sum_{j\in S_i}\omega_j q_j(X_i)$ reduces to $f_i(X_i)/q_i(X_i)$.
For $S_i=\{1,\ldots,N\}$ it gives the complete mixture ratio.  Only the selected
components change; the Gaussian density and derivative calculation is shared.
If every flow is the identity, $q_j=f_j$, both corrections are one.  These
limiting cases are useful implementation checks without implying that the
finite-particle answers must coincide for a nontrivial flow.

The score must differentiate the same mixture that supplied the value.
Let $D$ denote the total derivative along one parameter direction, including
the motion of the evaluated particle.  For a Gaussian component with mean
$\mu_j$ and covariance $\Sigma_j$, put
$r_{ij}=X_i-\mu_j$ and $u_{ij}=\Sigma_j^{-1}r_{ij}$.  Differentiating its
quadratic form and log determinant yields
\begin{equation}
 D\log\mathcal N(X_i;\mu_j,\Sigma_j)
 =-u_{ij}^{\top}(DX_i-D\mu_j)
 +\frac12u_{ij}^{\top}(D\Sigma_j)u_{ij}
 -\frac12\operatorname{tr}(\Sigma_j^{-1}D\Sigma_j).
 \label{eq:bf-hd-marginal-gaussian-tangent}
\end{equation}
For $\Sigma_j=B_jQB_j^\top$, all three terms in
$D\Sigma_j=(DB_j)QB_j^\top+B_j(DQ)B_j^\top+B_jQ(DB_j)^\top$
are required.  Omitting movement of the particle, the flow or the incoming
weights would give a partial derivative instead of the score claimed here.

For either component family $p_j=f_j$ or $p_j=q_j$, define its responsibilities
$\rho_{ij}^{p}=\omega_jp_j(X_i)/\sum_k\omega_kp_k(X_i)$.
Differentiating the sum before taking its logarithm gives
\begin{equation}
 D\log\bar p(X_i)
 =\sum_j\rho_{ij}^{p}
       \{D\log\omega_j+D\log p_j(X_i)\}.
 \label{eq:bf-hd-marginal-responsibility-score}
\end{equation}
Consequently $D\log a_i^{\rm mixture}$ is the sum of
$D\log\omega_i$, the total observation-likelihood derivative, and the
transition-mixture score minus the proposal-mixture score.  If
$Z_t=\sum_i a_i$ and $\widetilde\omega_i=a_i/Z_t$, the likelihood increment
and its derivative are $\log Z_t$ and
$D\log Z_t=\sum_i\widetilde\omega_iD\log a_i$.
The normalized-weight derivative passed to the next step is
$D\log\widetilde\omega_i=D\log a_i-D\log Z_t$.
Stable log-sum-exp reductions compute these expressions without exponentiating
large unnormalized log densities.

In the complete filtering step, the UKF predictions and per-ancestor affine
flows are formed first; the shared component evaluator then supplies either
ratio and its analytical total derivative.  The likelihood increment is
accumulated, weights are normalized, and the existing Contract E reset and
its analytical derivative propagate the cloud.  The capped marginal and
pairwise skewness/kurtosis corrections remain part of that reset.  Keeping
them does not guarantee that the empirical fourth moments, or the resulting
likelihood, are accurate: it preserves the stated algorithm while isolating
the effect of the importance correction.

A matched diagnostic on 6 October 2026 used one new length-50 dataset for each
model, exact prefixes of lengths 10, 20, 40 and 50, $N=1008$, and four paired
particle designs.  It used GPU FP64/XLA and fixed inherited controls, rather
than scope-specific tuning.  At $T=50$, predator--prey mean log likelihoods
were $-240.2772\pm0.0858$ with ancestor weights and
$-240.2759\pm0.0707$ with mixture weights, where the uncertainty is one
between-design standard error.  The author TT importance calculation gave
approximately $-240.2365$, and the $N=131072$ bootstrap reference gave
$-240.2401\pm0.0350$.  The small weighting difference does not establish a
ranking at this Monte Carlo precision.

SIR exposes a limitation that the predator--prey case does not.  At $T=10$,
the corresponding LEDH values were $-600.50\pm37.09$ and
$-579.66\pm25.88$, while the large bootstrap calculation gave
$-345.691\pm0.062$.  Even an ordinary $N=1008$ bootstrap calculation gave
$-346.007\pm0.348$.  A rank-20 author TT reference gave $-345.662$ with
conditional path standard error $0.004$; the independent rank-40 check was
still running at this documentation checkpoint.  These values show that
changing the importance correction alone does not repair this untuned SIR
configuration.  They do not identify whether its initial cloud, flow,
reset, or moment correction causes the discrepancy.  The completed values
and every score coordinate are preserved in the accompanying dated comparison
report; SIR score consistency and reference-rank checks are reported separately
from likelihood finiteness.

\section[Approximate targets and HMC]{Approximate Targets, Exact-Target Corrections, And HMC Admissibility}
\label{sec:bf-hd-fixed-branch-hmc}

For an unconstrained coordinate \(q\in\R^d\) and transform
\(\theta=\tau(q)\), the exact-target HMC potential is
\begin{equation}
  U(q)
  =
  -\ell_T(\tau(q))-
  \log \pi(\tau(q))-
  \log\left|\det D\tau(q)\right|.
  \label{eq:bf-hd-hmc-potential-live}
\end{equation}
If the filter supplies only \(\widehat\ell_T\), then ordinary HMC samples the
declared approximate posterior defined by \(\widehat\ell_T\).  Recovering the
exact posterior from an approximate or randomized likelihood would require a
separately defined exact-target extended-state or Metropolis correction scheme;
such a scheme is outside the scope of this chapter.

The HMC admissibility burden is therefore downstream of same-scalar discipline:
\begin{enumerate}[label=(\arabic*), leftmargin=0.28in]
  \item the filter must define the scalar likelihood or approximate likelihood
  before any sampler can be interpreted,
  \item the gradient used by HMC must be the gradient of the same scalar used in
  the Metropolis correction and posterior interpretation,
  \item transformed targets, including TT/KR-style maps, require the
  Jacobian-corrected target and valid support.
\end{enumerate}

\section{Exports To Validation And Promotion Chapters}
\label{sec:bf-hd-fixed-branch-exports}

The later validation chapter may import only the following consequences from the
present chapter:
\begin{itemize}[leftmargin=0.28in]
  \item the branch-valid finite-difference contract,
  \item the distinction between fixed structural ledgers and recomputed
  differentiable ledgers,
  \item the exact-target versus approximate-target HMC admissibility boundary,
  \item the rule that divergences, nonfinite gradients, and same-scalar mismatch
  are veto diagnostics rather than efficiency datapoints.
\end{itemize}

\section{Implementation Boundary}
\label{sec:bf-hd-hmc-boundary}

BayesFilter has no nonlinear HMC convergence result in this lane, no NeuTra
backend, no RMHMC backend, no tensor/KR transport-preconditioned HMC backend,
and no NAWM posterior recovery benchmark.  This chapter defines a shared branch,
scalar, and derivative contract.  It does not promote any sampler or backend as
production-ready for nonlinear state-space models.
